\begingroup
\def\C{\mathbb{C}}
\def\Tr{\operatorname{Tr}}
\def\cP{\mathcal P}
\section{Exact relation-window spectral transport on the original tau complex}
This calculation continues arithmetic determinant transport on its original
finite source over the absolute tau base. It computes the full spectra of
the two retained window operators, constructs an isometry intertwining
them in the actual source metric, and strengthens the signed determinant
bound. The source measures, their masses, the cyclic polynomial, the
Taylor unit, and the cohomology observation remain fixed.

\subsection{The original source and its four windows}

Use a full actual packet $h$ for $g=2\xi$, an integer $k\ge1$, and the
original cyclic polynomial $\chi=\chi_{h,k}$ of degree $q\ge1$.
Let $E=\C[S]/(\chi)$, $E_h=\C[s]/(h)$, and
$\upsilon_h=j_h(g/h)\in E_h^\times$.
The original tensor-sum observation is
\[
 \eta:E\hookrightarrow E_h^{\otimes k},\qquad
 \eta[P]=\upsilon_h^{\otimes k}P(A_h^{(k)})1.
 \tag{AW1}
\]
The original section $\sigma_h:E_h\hookrightarrow
Q=\mathscr B/\Theta V$ gives the separate injection
$\sigma_h^{\otimes k}\eta:E\hookrightarrow Q^{\otimes k}$.
Here $Q^{\otimes k}$ is the top cohomology of the original tensor
pushforward over $\mathbf F_{1,\tau}=\{\tau,1\}$; the one-factor
complex is $[V\oplus V\xrightarrow{\Theta(\phi-\widehat\psi)}
\mathscr B]$. The complete Tau Base and cyclic source proofs accompany
the preceding arithmetic determinant transport calculation (AT).

For $N\ge q-1$, retain the increasing-power spaces and maps
\[
 0\longrightarrow\cP_{N-q}\xrightarrow{B_N=\times\chi}
 \cP_N\xrightarrow{J_N}E\longrightarrow0,
 \qquad \cP_{-1}=0.
 \tag{AW2}
\]
The common source is $H=\cP_{2q}$, of dimension $n=2q+1$.
Its original coordinate on the integration line is $S=k/2+iu$.
Set $v_N(u)=(1,S(u),\ldots,S(u)^N)$ and retain the two densities
\[
 \begin{aligned}
 m_1(u)&=w_h^{*k}(u),&
 w_h(u)&=\frac{|(g/h)(1/2+iu)|^2}{2\pi},\\
 m_0(u)&=\frac{(\sqrt{2\pi})^k}{c_{k/4}}
              \frac{|\Gamma(k/4+iu/2)|^2}{2\pi},&
 c_a&=2^{1-2a}\Gamma(2a).
 \end{aligned}
 \tag{AW3}
\]
Their actual masses are $\mu_h^k$ and $(2\pi)^{k/2}$.
For the real parameter $0\le t\le1$, put
$m_t=(1-t)m_0+tm_1$, $M_N(t)=\int v_N^*v_Nm_t$, and
$M(t)=M_{2q}(t)$. The source density proofs give all finite moments
and positivity almost everywhere, so all these finite Grams are
positive definite. We use the original form
$\langle x,y\rangle_t=x^*M(t)y$ on $H$.

Let $P_N(t)$ be the orthogonal projection onto $\cP_N\subset H$,
and $Q_N(t)$ the projection onto $\chi\cP_{N-q}\subset H$.
Set $Q_{q-1}=0$. For the coefficient inclusion $I_N:\cP_N\to H$,
their fully typed coefficient formulas are
\[
 \begin{aligned}
 P_N&=I_NM_N^{-1}I_N^*M,\\
 Q_N&=I_NB_N(B_N^*M_NB_N)^{-1}B_N^*I_N^*M.
 \end{aligned}
 \tag{AW4}
\]
The zero relation domain at $N=q-1$ contributes the zero operator.
Let $G_N(t)$ be the minimum-lift Gram on $E$ from (AW2) and
$V_N(t)=\det G_N(t)$. Retain
\[
 \begin{aligned}
 \mathcal B(t)&=\log\frac{V_{q-1}(t)V_q(t)}
                              {V_{2q-1}(t)V_{2q}(t)},\\
 C(t)&=M(t)^{-1}(M(1)-M(0)),\\
 U(t)&=(Q_{2q-1}-Q_{q-1})+(Q_{2q}-Q_q),\\
 W(t)&=(P_{2q-1}-P_{q-1})+(P_{2q}-P_q).
 \end{aligned}
 \tag{AW5}
\]
The AT proof, included in full, establishes
$V_N=\det M_N/\det(B_N^*M_NB_N)$ with the original orientation
$(-1)^{q(N-q+1)}$, and
\[
 \Delta\mathcal B=\mathcal B(1)-\mathcal B(0)
       =\int_0^1\Tr((U(t)-W(t))C(t))\,dt.
 \tag{AW6}
\]
Thus the following computation evaluates the operators in the
already proved signed arithmetic correction.

\subsection{Full spectra and an explicit source isometry}

\begin{theorem}
In the original $M(t)$ metric, both $U(t)$ and $W(t)$ have spectrum
\[
 0\text{ with multiplicity }q,\qquad
 1\text{ with multiplicity }2,\qquad
 2\text{ with multiplicity }q-1.
 \tag{AW7}
\]
The last multiplicity is zero when $q=1$. There is an explicitly
constructed smooth $M(t)$-unitary $T(t):H\to H$ such that
$U(t)=T(t)W(t)T(t)^{-1}$.
\end{theorem}
\begin{proof}
Every subspace in either of the following flags is the literal
displayed polynomial subspace, with its original coefficient inclusion.
For the relation flag use
\[
 F_0=\chi\cP_0\subset F_1=\chi\cP_{q-1}
            \subset F_2=\chi\cP_q\subset H.
 \tag{AW8}
\]
Their dimensions are $1,q,q+1$, since multiplication by $\chi$ is
injective. When $q=1$, $F_0=F_1$ is retained as that equality.
Projections of nested subspaces commute and their products equal
the projection of the smaller subspace: each larger space is the
orthogonal direct sum of the smaller space and its orthogonal
complement within the larger space. On the original orthogonal
decomposition
$F_0\oplus(F_1\cap F_0^\perp)\oplus
(F_2\cap F_1^\perp)\oplus F_2^\perp$, the operator
$U=Q_{F_1}+Q_{F_2}-Q_{F_0}$ acts respectively by $1,2,1,0$.
The dimensions are respectively $1,q-1,1,q$.

For $W$ use
$\cP_{q-1}\subset\cP_q\subset\cP_{2q-1}\subset\cP_{2q}=H$.
On the four corresponding orthogonal pieces, $W$ acts by
$0,1,2,1$, with dimensions $q,1,q-1,1$. This proves (AW7),
including the repeated middle subspaces at $q=1$.

For an explicit intertwiner retain every Gram factor as follows.
For any specified ordered independent list $z_0,\ldots,z_d$,
apply the recursion
\[
 y_j=z_j-\sum_{i<j}e_i\langle e_i,z_j\rangle_t,
 \qquad a_j=(y_j^*M(t)y_j)^{1/2}>0,
 \qquad e_j=y_j/a_j.
 \tag{AW9}
\]
The positive real square root fixes its phase. No source density or
mass is changed; each $a_j$ is the norm in the original form.
Independence proves $y_j\ne0$ at every step, so this recursion is
defined and smooth on the entire parameter interval.

First use $z_j=S^j$, $0\le j\le2q$, obtaining $p_0,\ldots,p_{2q}$.
Next use $z_j=\chi S^j$, $0\le j\le q$, obtaining
$r_0,\ldots,r_q$. These lie in the successive relation spaces
in (AW8). Finally let $R_{2q}(t):E\to H$ be the original minimum
section from AT6 and apply (AW9) to
$R_{2q}(t)[1],\ldots,R_{2q}(t)[S^{q-1}]$, obtaining
$c_0,\ldots,c_{q-1}$. The relation orthogonality in AT6 puts this
entire list in $F_2^\perp$. Since $J_{2q}R_{2q}=I_E$, the list
is independent and spans that $q$-dimensional complement.

Define $T$ on the first orthonormal basis by
\[
 \begin{aligned}
 Tp_j&=c_j&& (0\le j\le q-1),\\
 Tp_q&=r_0,&\quad Tp_{2q}&=r_q,\\
 Tp_{q+j}&=r_j&& (1\le j\le q-1).
 \end{aligned}
 \tag{AW10}
\]
The middle range is empty for $q=1$, with $p_q$ and $p_{2q}$
still distinct. The target list is an orthonormal basis. Thus this
linear map is invertible and $T^*MT=M$. The spectral actions
already calculated agree on each indicated basis vector, proving
$UT=TW$. The recursion, the smooth minimum section and the positive
norms prove the asserted smoothness. This completes the construction.
\end{proof}

The exact isometry gives a commutator form of the signed integrand:
\[
 \Tr((U-W)C)=\Tr\bigl(WT^{-1}[C,T]\bigr),
 \qquad [C,T]=CT-TC.
 \tag{AW11}
\]
Indeed cyclicity moves $T$ in $\Tr(TWT^{-1}C)$ to the right,
and subtraction then gives the displayed product.
This $T$ is an isometry of the original source; its cohomology
observation is computed, rather than silently assigned an
intertwining property. The unchanged arithmetic observation
$\mathcal O=\sigma_h^{\otimes k}\eta J_{2q}:H\to Q^{\otimes k}$
satisfies
\[
 \mathcal O T=\sigma_h^{\otimes k}\eta J_{2q}T,
 \quad \mathcal O T-\mathcal O
       =\sigma_h^{\otimes k}\eta J_{2q}(T-I),
 \quad \ker(\mathcal O T)=T^{-1}F_2.
 \tag{AW12}
\]
The last identity follows from both injections in (AW1) and
$\ker J_{2q}=F_2$. By (AW10), $T^{-1}F_2$ is exactly
$\operatorname{span}\{p_q,\ldots,p_{2q}\}
=\cP_{q-1}^{\perp}$ in $H$.
At any fixed nonempty support label $A$, lift every displayed
linear map by $(A,x)\mapsto(A,fx)$ and send $\tau$ to $\tau$.
Consequently the supported-zero inverse image of the lifted
$\mathcal O T$ at $A$ is the full set
$\{A\}\times\cP_{q-1}^{\perp}$. The Taylor unit in (AW1)
and the defect in (AW12) remain present.

\subsection{The spectral bound with every eigenvalue retained}

Let $b_1\le\cdots\le b_n$ be all eigenvalues, counted with
multiplicity, of $D=M(0)^{-1}M(1)$ on its original $M(0)$
Hilbert space. All are positive. Define
\[
 \mathcal L_q(D)=\log\frac{b_{q+2}}{b_q}
       +2\sum_{j=1}^{q-1}\log\frac{b_{2q+2-j}}{b_j},
 \qquad \kappa(D)=b_n/b_1.
 \tag{AW13}
\]
For $q=1$ the sum is empty and $\mathcal L_1(D)=\log(b_3/b_1)$.


Retain also the actual polynomial-space overlap. At each $t$, let
$s(t)=\dim(\operatorname{ran}U\cap\operatorname{ran}W)$ and
$d(t)=c_n(t)-c_1(t)$. With the actual $c_j(t)$ of (AW17), set
\[
 \begin{aligned}
 S_{\rm AW}(t)&=c_{q+2}-c_q+
          2\sum_{j=1}^{q-1}(c_{2q+2-j}-c_j),\\
 S_{\rm SP}(t)&=\frac{d(t)}2\min\left\{4q-2s(t),
       \sqrt{(2q+1)(8q-4-2\Tr(UW))}\right\},\\
 \mathcal J_{h,k}&=\int_0^1\min\{S_{\rm AW}(t),S_{\rm SP}(t)\}\,dt.
 \end{aligned}
 \tag{AW13a}
\]
The overlap refinement from SP8--SP10 is transported to this
earlier calculation by the exact arrows
$x=t$, $y=u$, $B_N^{\rm SP}=I_NB_N^{\rm AW}$.
Indeed $M_N=I_N^*MI_N$ gives
$(B_N^{\rm SP})^*MB_N^{\rm SP}=B_N^*M_NB_N$, so its projectors
are precisely (AW4), and $(\mathcal R^{\rm SP},\mathcal P^{\rm SP},
T^{\rm SP})=(U,W,C)$. The original Gamma densities agree because
$r_{1/4}^{*k}=c_{1/4}^{\,k}r_{k/4}/c_{k/4}$ with
$c_{1/4}=\sqrt{2\pi}$, exactly (AW3). The arithmetic densities
already coincide as $w_h^{*k}$. Thus this comparison preserves
the moments, masses, polynomial coefficients and all four endpoint signs.



The exact trace norm and the centered Hilbert--Schmidt estimate
are also retained before taking a minimum. On this same source define
\[
 \begin{aligned}
 S_{\rm tr}(t)&=\frac{d(t)}{2}\|U(t)-W(t)\|_1,\\
 \Sigma(t)&=\operatorname{Tr}(C(t)^2)
       -\frac{(\operatorname{Tr}C(t))^2}{2q+1},\\
 S_{\rm HS}(t)&=\sqrt{\Sigma(t)\operatorname{Tr}((U(t)-W(t))^2)}.
 \end{aligned}
 \tag{AW13c}
\]
The trace-norm inequality follows by subtracting the midpoint
of the extreme eigenvalues from the metric derivative, since the
trace of the signed projection difference is zero. In an orthonormal
eigenbasis of that difference, every diagonal derivative entry is
bounded in modulus by half the spectral diameter; summing with
the absolute eigenvalues proves the assertion.
For the other bound subtract the trace mean
$\operatorname{Tr}(C(t))I/(2q+1)$ instead.
The square of its Hilbert--Schmidt norm in the unchanged source
metric is exactly $\Sigma(t)$. Hilbert--Schmidt Cauchy--Schwarz
with the signed difference proves that its trace pairing is bounded
by $S_{\rm HS}(t)$. Positivity of both factors under the square
root follows from their self-adjointness. Thus both functions bound
the same absolute signed derivative on general positive source forms;
the angle term below still uses the specified moment measures.
These are RMT10a--10b and ISM25--26 through the already proved
original coefficient and projection maps.

The original-moment angle calculation supplies a further simultaneous
bound at this same claim site. Put $a_q=2q-1$ and
\[
 \begin{aligned}
 A_{h,k}(t)&=a_q\sqrt{1-e^{-\mathcal B(t)/a_q}}\,d(t),\\
 \mathcal J_{h,k}^{(5)}
   &=\int_0^1\min\{S_{\rm AW}(t),S_{\rm tr}(t),S_{\rm SP}(t),S_{\rm HS}(t),A_{h,k}(t)\}\,dt,\\
 \mathcal I_{h,k}
   &=\int_0^1\min\left\{d(t),
   \frac{S_{\rm AW}(t)}{a_q\sqrt{1-e^{-\mathcal B(t)/a_q}}},
   \frac{S_{\rm SP}(t)}{a_q\sqrt{1-e^{-\mathcal B(t)/a_q}}},\right.
 \\
 &\hspace{2em}\left.   \frac{S_{\rm tr}(t)}{a_q\sqrt{1-e^{-\mathcal B(t)/a_q}}},
   \frac{S_{\rm HS}(t)}{a_q\sqrt{1-e^{-\mathcal B(t)/a_q}}}\right\}\,dt .
 \end{aligned}
 \tag{AW13b}
\]
These specialize RMT1--RMT18 through $a_{\rm RMT}=a_q$, $x=t$, $C_x=C(t)$ and
$(U_x,W_x)=(U(t),W(t))$. The density and relation arrows were
proved in (AW3)--(AW4) and (AW13a), so every source moment and original
cohomology map agrees. This application uses the actual moment
measures (AW3). It does not impose moment structure on an
arbitrary positive coefficient Gram.

\begin{theorem}
On the unchanged arithmetic source,
\[
 |\Delta\mathcal B|\le\mathcal J_{h,k}^{(5)}\le\mathcal J_{h,k}\le\mathcal L_q(D)
                  \le(2q-1)\log\kappa(D).
 \tag{AW14}
\]
The previously proved coefficient $2q-1$ and full relative-spectrum
bound are retained. The pointwise minimum additionally preserves the
actual polynomial-space overlap, and is at most $\int_0^1S_{\rm SP}(t)\,dt$.
\end{theorem}
\begin{proof}
We first prove the needed finite spectral trace inequality.
Let $C$ be self-adjoint in a finite Hilbert space with eigenvalues
$c_1\le\cdots\le c_n$. If $P$ is an orthogonal projection of
rank $r$, let $e_j$ be an orthonormal eigenbasis of $C$ and
$d_j=\langle e_j,Pe_j\rangle$. Then $0\le d_j\le1$ and
$\sum_jd_j=r$, since $d_j=\|Pe_j\|^2$ and $\Tr P=r$.
Thus $\Tr(PC)=\sum_jc_jd_j$.
Its upper bound by $\sum_{j=n-r+1}^n c_j$ follows by subtracting
the two expressions: the mass $\sum_{j\le n-r}d_j$ equals
$\sum_{j>n-r}(1-d_j)$, and every coefficient in the latter
group is at least every coefficient in the former. Pairing the
same nonnegative mass, or bounding both groups by $c_{n-r+1}$,
proves the inequality. Applying this argument to $-C$ proves
the lower bound $\sum_{j=1}^r c_j$. This proof includes repeated
eigenvalues and the ranks zero and $n$.

For an operator $Z$ with spectrum (AW7), its spectral projection
$A_Z$ onto eigenvalues $1,2$ has rank $q+1$, and its spectral
projection $B_Z$ onto eigenvalue $2$ has rank $q-1$.
The exact decomposition is $Z=A_Z+B_Z$. Applying the proved
projection bounds gives
\[
 \sum_{j=1}^{q+1}c_j+\sum_{j=1}^{q-1}c_j
 \ \le\ \Tr(ZC)\ \le\
 \sum_{j=q+1}^{2q+1}c_j+\sum_{j=q+3}^{2q+1}c_j.
 \tag{AW15}
\]
Each bound is attained for an abstract operator of this spectrum
by assigning its eigenspaces to the appropriate eigenvectors
of $C$. This attainment statement is only about the displayed
spectral constraint; it assigns no freely chosen orientation
to the actual polynomial relation spaces.
Both $U$ and $W$ satisfy (AW15). Subtract their bounds and cancel
the common $c_{q+1}$ term only after writing both complete sums:
\[
 |\Tr((U-W)C)|
 \le c_{q+2}-c_q
       +2\sum_{j=1}^{q-1}(c_{2q+2-j}-c_j).
 \tag{AW16}
\]
No mutual commutation of $U,W,C$ is used.

The range dimensions in (AW7) give $s(t)\ge1$. Let $E_L$ be the
orthogonal projector onto their intersection. Each $Z=U,W$ is
at least its range projection, since its nonzero eigenvalues are
$1,2$, and therefore $Z-E_L$ is positive of trace $2q-s(t)$.
Writing $A=U-W=(U-E_L)-(W-E_L)$ proves
\[
 \|A\|_1\le4q-2s(t),\quad
 \Tr A=0,\quad
 \Tr A^2=8q-4-2\Tr(UW),\quad
 \|A\|_1\le\sqrt{(2q+1)\Tr A^2}.
 \tag{AW16a}
\]
The square identity uses $\Tr U^2=\Tr W^2=4q-2$ and cyclicity.
The last inequality is Cauchy--Schwarz on every real eigenvalue of
the self-adjoint $A$. Subtracting $(c_1+c_n)I/2$ from $C$
does not change its trace pairing with $A$. In an orthonormal
eigenbasis of $A$, that pairing is bounded in absolute value by
$d(t)\|A\|_1/2$. This proves
$|\Tr((U-W)C)|\le S_{\rm SP}(t)$ in addition to (AW16),
and hence bounds it by their pointwise minimum. Smooth positive
Grams give bounded continuous eigenvalues and overlaps;
$s(t)$ is Borel because its rank strata are specified by matrix
minors. Therefore the minimum is integrable. The signed identity
(AW6) now proves $|\Delta\mathcal B|\le\mathcal J_{h,k}$ and
$\mathcal J_{h,k}\le\int S_{\rm SP}$.
For the original moment path, the exact original-source identification
above also applies RMT11--RMT12: its full crossed angle lists give
$|\mathcal B'|\le A_{h,k}$, and $\mathcal B>0$ follows because
otherwise $1\perp\chi\mathcal P_q$ would force
$\int|\chi|^2m_t=0$ by the reflected polynomial
$\chi^{\#_k}(S)=\sum_j\overline{\chi_j}(k-S)^j$.
Taking the pointwise minimum of this additional proved inequality
with (AW16), (AW16a), and both complete bounds in (AW13c) gives
$|\Delta\mathcal B|\le\mathcal J_{h,k}^{(5)}$.
The complete proof of the original angle estimate and its exact
chain-rule continuation is retained in EP and RMT; it uses
AW's earlier projection identities, not the final claim (AW14).


For the actual interpolated metric,
$C(t)=((1-t)I+tD)^{-1}(D-I)$, so its eigenvalues in increasing
order are
\[
 c_j(t)=\frac{b_j-1}{1-t+tb_j}.
 \tag{AW17}
\]
The denominator is positive and the derivative with respect
to $b_j$ is $(1-t+tb_j)^{-2}>0$. Thus the ordering is valid
for the entire closed interval. Each $c_j(t)$ is the derivative
of $\log(1-t+tb_j)$ and has integral $\log b_j$.
Integrating (AW16) proves $\mathcal J_{h,k}\le\mathcal L_q(D)$;
the preceding pointwise-minimum argument proves the first inequality
of (AW14). Every ratio in (AW13) lies
between $1$ and $b_n/b_1$, so retaining its coefficient gives
$1+2(q-1)=2q-1$ times $\log\kappa(D)$ for the second bound.
This proves the theorem for all $q\ge1$.
\end{proof}

All bounds retain the literal source masses. If $M(1)=cM(0)$,
every $b_j=c$, so (AW13) is zero. Directly each original
quotient determinant is multiplied at time $t$ by
$(1-t+tc)^q$, and the four displayed factors in (AW5) cancel.
For a fixed invertible source-coordinate arrow $F$, the transported
forms are $M'=F^*MF$, all subspaces have their $F^{-1}$ images,
and $D'=F^{-1}DF$. Hence the eigenvalues, (AW6) and (AW13)
are unchanged under this explicit map.


For every integer $\ell_N\ge0$ the constructed signed approximation
is applied on these exact endpoint forms as well. Here $\mathsf M_z=M(t)|_{t=z}$,
$\mathsf C_z=C(t)|_{t=z}$ and
$\mathsf A_z=(U(t)-W(t))|_{t=z}$, with $\mathsf D=D$.
The exact RMT/ISM source arrow is the identity,
and $\mathcal B$ is the function (AW5).
Write $\mathsf D=(\mathsf M_0)^{-1}\mathsf M_1$, with its complete
ordered positive eigenvalues $b_1,\ldots,b_{2q+1}$, and let
$\mathsf A_z$ be this same signed projection difference. Define
\[
 \begin{aligned}
 \alpha_z^\circ&=\frac{2(1-z)+z(b_1+b_{2q+1})}{2},&
 H_z^\circ&=I-\frac{(1-z)I+z\mathsf D}{\alpha_z^\circ},\\
 C_z^{[\ell_N]}&=(\alpha_z^\circ)^{-1}
       \sum_{r=0}^{\ell_N}(H_z^\circ)^r(\mathsf D-I),&
 R_z^{[\ell_N]}&=(H_z^\circ)^{\ell_N+1}\mathsf C_z,\\
 E_z^{[\ell_N]}&=R_z^{[\ell_N]}
          -\frac{\operatorname{Tr}R_z^{[\ell_N]}}{2q+1}I .
 \end{aligned}
 \tag{AW17aa}
\]
The symbols $\alpha_z^\circ,H_z^\circ$ preserve every original
source parameter and relation map; $\ell_N$ is only the finite
series truncation. Multiplication of the finite geometric sum
gives $\mathsf C_z-C_z^{[\ell_N]}=R_z^{[\ell_N]}$.
All factors are real functions of the positive relative
operator $\mathsf D$, hence self-adjoint in the original
$\mathsf M_z=\mathsf M_0((1-z)I+z\mathsf D)$ metric.
Since $\operatorname{Tr}\mathsf A_z=0$, the derivative of the same
four-volume function is its signed center
$\operatorname{Tr}(C_z^{[\ell_N]}\mathsf A_z)$ plus
$\operatorname{Tr}(E_z^{[\ell_N]}\mathsf A_z)$.
Consequently
\[
 \begin{aligned}
 j_{\ell_N}&=\int_0^1
       \operatorname{Tr}(C_z^{[\ell_N]}\mathsf A_z)\,dz,\\
 e_{\ell_N}&=\int_0^1
       \sqrt{\operatorname{Tr}((E_z^{[\ell_N]})^2)
                     \operatorname{Tr}(\mathsf A_z^2)}\,dz,\\
 j_{\ell_N}-e_{\ell_N}
 &\le\mathcal B(1)-\mathcal B(0)\le j_{\ell_N}+e_{\ell_N}.
 \end{aligned}
 \tag{AW17ab}
\]
For this block $\mathcal B$ denotes precisely the four-volume
function identified above. The signed inequalities follow by
Hilbert--Schmidt Cauchy--Schwarz in that unchanged metric and
integration of the exact derivative.
The full residual calculation ISM29--37, equivalently RMT19--21,
uses
$\vartheta=(b_{2q+1}-b_1)/(b_{2q+1}+b_1)<1$ and
$K_{\mathsf D}=\sum_j(|b_j-1|/\min\{1,b_j\})^2$ to prove
\[
 0\le e_{\ell_N}\le
    \vartheta^{\ell_N+1}\sqrt{K_{\mathsf D}(8q-6)}
       \longrightarrow0.
 \tag{AW17ac}
\]
Indeed every eigenvalue of $H_z^\circ$ has modulus at most
$\vartheta$, centering the residual subtracts a nonnegative square
from its squared Hilbert--Schmidt norm, and
$\operatorname{Tr}(\mathsf A_z^2)\le8q-6$ by the full flag spectra
and the common range. This proves the bound with all eigenvalue
multiplicities retained. The displayed integrals are exact; a
numerical evaluation must retain its own integration error.
Intersecting this proved signed interval with the already proved
absolute and nonlinear intervals supplies the additional endpoint
entries below, for every specified $\ell_N$.

The original quartet Gamma theorem, when applied to the actual
full quartet packet with displacement $\delta>0$ and common
order $m$, retains $q=[1+k(m-1)](k+1)^2$ and gives
\[
 \mathcal B^{\rm ar}_{h,k}
 \ge4q\log\!\left(\frac{\delta k}{2\sqrt5}\right)
             +\max\{-\mathcal J_{h,k}^{(5)},j_{\ell_N}-e_{\ell_N}\}.
 \tag{AW18}
\]
This substitutes the complete supplied Gamma theorem into the
lower branches of (AW14) and the constructed signed inequality (AW17ab),
then takes their maximum. It retains the stronger
actual overlap bound, and also implies the earlier lower bounds with
$-\mathcal L_q(D)$ and $-(2q-1)\log\kappa(D)$ respectively. It proves no
uniform estimate for the eigenvalues $b_j$, and asserts no
existence of an off-line packet. The new computed objects are
the exact window spectra, their source isometry, its arithmetic
observation and commutator, and the strengthened finite bound.

\section[The one-dimensional quotient: exact angle and nonlinear transport]{The one-dimensional quotient:\\exact angle and nonlinear transport}

When $q=1$, the same four endpoints are $0,1,1,2$. Their
middle factors cancel in the specified ratio, leaving
$\mathcal B(t)=\log(V_0(t)/V_2(t))$ with both original masses
still present. Write $\chi=S-a$. The actual operators now are
$U=Q_2$ and $W=I-P_0$. If $\mathcal N=(\chi\cP_1)^\perp$
and $P_{\mathcal N}=I-Q_2$, then
$U-W=P_0-P_{\mathcal N}$. Both projections have rank one.

Let $e=1\in H$, $z=R_2(t)[1]$, and retain their actual norms
$\langle e,e\rangle_t=V_0(t)$ and
$\langle z,z\rangle_t=V_2(t)$. The source splitting gives
$z=P_{\mathcal N}e$, since $J_2e=1$. Thus
$\langle e,z\rangle_t=V_2(t)$, and the angle between the
two lines has the exact value
\[
 \cos^2\theta(t)=\frac{|\langle e,z\rangle_t|^2}
                         {\langle e,e\rangle_t\langle z,z\rangle_t}
       =\frac{V_2(t)}{V_0(t)}=e^{-\mathcal B(t)}.
 \tag{AW19}
\]
Both Grams are positive, so the cosine is positive. It is
strictly less than one for the actual measure (AW3).
Otherwise $e\in\mathcal N$, implying
$\int\chi(S)m_t=0$ and $\int S\chi(S)m_t=0$.
On the unchanged line $\overline S=k-S$, and consequently
$|\chi(S)|^2=((k-\overline a)-S)\chi(S)$.
The two equalities would force $\int|\chi(S)|^2m_t=0$,
contrary to the positive Gram of this nonzero polynomial.
Hence $0<\theta(t)<\pi/2$ and $\mathcal B(t)>0$ on the
whole closed interval.

For completeness, take the unit vector $p=e/\sqrt{V_0}$
and write $z/\sqrt{V_2}=\cos\theta\,p+\sin\theta\,r$,
where $r$ is the resulting unit vector orthogonal to $p$.
The phase is already fixed since $\langle e,z\rangle_t=V_2>0$.
In this orthonormal pair the exact difference of projections is
\[
 \left.(P_0-P_{\mathcal N})\right|_{\operatorname{span}(p,r)}
 =\begin{pmatrix}
   \sin^2\theta&-\sin\theta\cos\theta\\
   -\sin\theta\cos\theta&-\sin^2\theta
  \end{pmatrix}.
 \tag{AW20}
\]
Its trace is zero and its determinant is $-\sin^2\theta$.
It vanishes on the remaining one-dimensional orthogonal
complement. Thus its full eigenvalues are
$\sin\theta,0,-\sin\theta$.
Applying the proved rank-one spectral trace bounds to its
positive and negative eigenprojections yields
\[
 |\dot{\mathcal B}(t)|
 \le\sqrt{1-e^{-\mathcal B(t)}}\,(c_3(t)-c_1(t)).
 \tag{AW21}
\]
This retains the actual alignment of the two lines.

The function $\Phi(B)=2\operatorname{arcosh}(e^{B/2})$
for $B>0$ has derivative
$\Phi'(B)=1/\sqrt{1-e^{-B}}$, by differentiating the
inverse of the hyperbolic cosine on its positive branch.
The preceding strict positivity on the compact interval
justifies its smooth composition with $\mathcal B(t)$.
Divide (AW21) by its positive square-root factor. The four
independent generic inequalities (AW16), (AW16a), and the trace-norm
and Hilbert--Schmidt inequalities (AW13c) can all be divided by
that same factor because $q=1$ gives $a_q=1$ and $\mathcal B>0$.
Their pointwise minimum integrates to $\mathcal I_{h,k}$ in
(AW13b); integration of (AW17) bounds it by $\log(b_3/b_1)$. The exact nonlinear endpoint bound is
\[
 \left|2\operatorname{arcosh}(e^{\mathcal B(1)/2})
       -2\operatorname{arcosh}(e^{\mathcal B(0)/2})\right|
       \le\mathcal I_{h,k}\le\log(b_3/b_1).
 \tag{AW22}
\]
Combining (AW22) with the signed interval (AW17ab), with the original values
$A_0=\operatorname{arcosh}(e^{\mathcal B(0)/2})$
and $L=\mathcal I_{h,k}/2\le\tfrac12\log(b_3/b_1)$, monotonicity of
$2\log\cosh A$ for $A\ge0$ gives
\[
 \begin{aligned}
 \mathcal B(1)&\ge
 \max\{2\log\cosh(\max\{0,A_0-L\}),
          \mathcal B(0)+j_{\ell_N}-e_{\ell_N}\},\\
 \mathcal B(1)&\le
 \min\{2\log\cosh(A_0+L),
          \mathcal B(0)+j_{\ell_N}+e_{\ell_N}\}.
 \end{aligned}
 \tag{AW23}
\]
This last lower bound can be zero although the actual
$\mathcal B(1)$ is strictly positive, as proved above.
Every step concerns the original $q=1$ quotient and its
degree-two source. No extension of that dimension restriction
to a quartet quotient is asserted.

\subsection{Complete proof dependencies}

The complete AT calculation accompanies this note, together with its
29 complete source dependencies and independent proof review. Those
sources prove the original pushforward, tensor signs, full Taylor-unit
injections, theta primitive, positive arithmetic and Gamma moments,
and the quartet application. Every new assertion (AW7)--(AW23) is
proved above on those unchanged finite objects. A separate independent
review supplies the spectral calculation and its edge-case checks.

\endgroup
